%% kapa: exam=1 round=29 year=2018 period=1 form=A subject=civil_law subject_ko=민법 seq=40 number=40
%% source: Q-Net 감정평가사 기출문제 (cst00309) artlSeq=5202808 file=extracted/1008/5202808/2018년 제29회 감정평가사 1차시험 기출문제/2018년 제29회 감정평가사 1차시험 기출문제/1교시 A형.pdf
%% conversion: kordoc 4.24.51 -> tools/gen1.py (auto); review=proofread-sonnet; flags=none
\begin{question}{40}
동산질권에 관한 설명으로 옳지 않은 것은? (다툼이 있으면 판례에 따름)
\begin{choices}
  \item 동산질권도 선의취득의 대상이 될 수 있다.
  \item 질권설정자는 채무변제기 후의 계약으로 질권자에게 변제에 갈음하여 질물의 소유권을 취득하게 할 것을 약정하지 못한다.
  \item 수개의 채권을 담보하기 위하여 동일한 동산에 수개의 질권을 설정한 때에는 그 순위는 설정의 선후에 의한다.
  \item 다른 약정이 없는 한 질권은 원본, 이자, 위약금, 질권실행의 비용, 질물보존의 비용 및 채무불이행 또는 질물의 하자로 인한 손해배상의 채권을 담보한다.
  \item 정당한 이유가 있는 경우 질권자는 간이변제충당을 법원에 청구할 수 있고, 이 때 질권자는 미리 채무자 및 질권설정자에게 통지하여야 한다.
\end{choices}
\end{question}
\answer{2}
